% Part: lambda-calculus
% Chapter: introduction
% Section: lambda-computable

\documentclass[../../../include/open-logic-section]{subfiles}

\begin{document}

\olfileid[hi]{lam}{int}{cmp}
\olsection{\usetoken{S}{lambda definable} फलन संगणनीय हैं}

\begin{thm}
\ollabel{thm:lambda-computable}
यदि कोई आंशिक फलन $f$ किसी लैम्ब्डा पद द्वारा !!{lambda defined} है, तो वह
संगणनीय है।
\end{thm}

\begin{proof}
मान लें फलन~$f$, किसी लैम्ब्डा पद~$X$ द्वारा !!{lambda defined} है।
~$f$ को संगणित करने की एक अनौपचारिक प्रक्रिया बताएँ। निवेश $m_0$,
\dots,~$m_{n-1}$ पर पद $X \num m_0 \ldots \num m_{n-1}$ लिखें। एक वृक्ष
बनाएँ: पहले मूल पद के सभी एक-चरणीय अपचयन लिखें; उनके नीचे उन सबके सभी
एक-चरणीय अपचयन (अर्थात् मूल पद के दो-चरणीय अपचयन) लिखें; और इसी प्रकार आगे
बढ़ते रहें। यदि कभी किसी अंक तक पहुँचें, तो उत्तर के रूप में वही लौटाएँ;
अन्यथा फलन अपरिभाषित है।

चर्च की प्रस्थापना का आह्वान हमें बताता है कि यह फलन संगणनीय है। प्रमेय को
सिद्ध करने का बेहतर उपाय इस खोज-प्रक्रिया का पुनरावर्ती वर्णन देना होगा।
उदाहरणार्थ, हम आदिम पुनरावर्ती फलनों और संबंधों का एक अनुक्रम परिभाषित कर
सकते हैं: ``$\fn{IsASubterm}$,'' ``$\fn{Substitute}$,''
``$\fn{ReducesToInOneStep}$,'' ``$\fn{ReductionSequence}$,''
``$\fn{Numeral}$,'' आदि। तब $f(m_0, \dots, m_{n-1})$ की संगणना की आंशिक
पुनरावर्ती प्रक्रिया ऐसे एक-चरणीय अपचयनों के अनुक्रम की खोज करना है जो $X
\num{m_0} \dots \num{m_{n-1}}$ से आरंभ होकर किसी अंक पर समाप्त हो, और उस
अंक के अनुरूप संख्या लौटाना है। विवरण लंबे और नीरस हैं, किंतु अन्यथा
नियमित हैं।
\end{proof}

\end{document}
